\section{Multiplicity and the proposed second-eigenvalue gate}
\label{sec:spectral-multiplicity-gate}

\begin{proposition}[AC-S34: a multiple zero forces two small even modes]
\label{AC-S34}
Assume $A$.  If a nontrivial zeta zero at ordinate
$\gamma\ne0$ has multiplicity at least two, then,
for every $\varepsilon>0$,
\[
 a_2(\log\lambda)\leq
        C_{\gamma,\varepsilon}\lambda^{2+\varepsilon}(1-\nu_4)
 \leq C'_{\gamma,\varepsilon}
                \lambda^{11+\varepsilon}e^{-4\pi\lambda^2}.
                                                        \tag{SM.1}
\]
The second eigenvalue in the even logarithmic sector
satisfies the same estimate.  Consequently the additional
lower bound (SE.14) in AC-S33, together with $A$, would
imply simplicity of every nontrivial zeta zero.
This is a limitation of that sufficient gap criterion;
it is not an assertion that a multiple zero exists,
or that every possible lowest-line convergence theorem
would require simplicity.
\end{proposition}
\begin{proof}
We may take $\gamma>0$, since the conjugate ordinate
has the same multiplicity.  Let $k_\lambda^+$ denote
the even part of the original trial of AC-S30,
and write $\delta=1-\nu_4$, $L=\log\lambda$.

\proofstep{1. The two jet errors have exponentially small size.}
The full, uncut theta function
$F_\lambda(x)=E\psi_\lambda(e^x)$ has Mellin transform
$\zeta(s)\int\psi_\lambda(u)u^{s-1}du$ on
$\operatorname{Re}s>0$, as proved in AC-S30.
The raw Mellin factor is holomorphic at the critical
zero.  Therefore the even part of $F_\lambda$ has
Fourier value and first derivative equal to zero
at both $\gamma$ and $-\gamma$.
The errors for the cut function $k_\lambda^+$ are
exactly the corresponding Fourier jets of its two
omitted tails.

The endpoint estimate of AC-S31 holds for every
fixed logarithmic moment.  The multiplier argument
of (SR.16) applies to each fixed derivative of
$\zeta$: the same fixed smaller disk and Cauchy's
formula give an arbitrarily small power bound for
that derivative.  Differentiating the multiplier
twice therefore shows, for the reflected and shifted
tail $R$ of (SR.19),
\[
 \|(1+t^2)R\|_2
             \leq C_\sigma\lambda^{1+4\sigma}\sqrt\delta
                       \quad\text{for every }\sigma>0
                       \text{ sufficiently small}. \tag{SM.2}
\]
The raw moments needed are at most two; the sharp
half-line cutoff commutes with those moments, and
the scalar exponential tail has all these moments.
Cauchy--Schwarz now gives the same bound for
$\int(1+|t|)|R(t)|dt$.
Translating a tail by $L$ introduces at most a factor
$1+L$ in its first Fourier derivative.
Thus the real two-vector of errors satisfies
\[
 e_\lambda=
     \bigl(\widehat{k_\lambda^+}(\gamma),
                    (\widehat{k_\lambda^+})'(\gamma)\bigr),
 \qquad |e_\lambda|\leq
       C_{\gamma,\sigma}(1+L)
                         \lambda^{1+4\sigma}\sqrt\delta.
                                                        \tag{SM.3}
\]
The Fourier convention here is $e^{-itx}$.

\proofstep{2. Exact jet correction in a fixed physical interval.}
Fix $a>0$ and a real nonnegative even smooth function
$\theta$ supported in $(-a,a)$ and positive near zero.
Put
\[
 v_0(x)=\cos(\gamma x),\quad
 v_1(x)=-x\sin(\gamma x),\quad
 b_j(x)=\theta(x)v_j(x),\quad
 M_{ij}=\int\theta(x)v_i(x)v_j(x)dx.              \tag{SM.4}
\]
These are even real compact profiles.
The matrix $M$ is positive definite: the two analytic
functions $v_0,v_1$ are linearly independent on an
interval around zero, as is seen first at zero and
then from $\gamma\ne0$.
Moreover $M_{0j}=\widehat b_j(\gamma)$ and
$M_{1j}=\widehat b_j'(\gamma)$.
For $L>a$ define
\[
        j_\lambda=k_\lambda^+
                    -(b_0,b_1)M^{-1}e_\lambda.     \tag{SM.5}
\]
Its Fourier transform has exact double zeros at
$\gamma$ and $-\gamma$; the latter conditions follow
by evenness.  The correction has fixed compact
support and fixed finite Weil form norm.
Positivity of the form, (SM.3), and AC-S32 give
\[
 q_L(j_\lambda)
       \leq C_{\gamma,\varepsilon}
                                  \lambda^{2+\varepsilon}\delta
                      \quad\text{for every }\varepsilon>0.
                                                        \tag{SM.6}
\]
Here first choose $\sigma$ small and then absorb
$(1+L)^2$ into the remaining arbitrary power of
$\lambda$.  By the scalar leakage asymptotic already
used in AC-S30, the correction is exponentially
small in $L^2$.  Hence
\[
                     j_\lambda\longrightarrow K
            \quad\text{with error }O(\lambda^{-1/2})\text{ in }L^2.
                                                        \tag{SM.7}
\]

\proofstep{3. Compactly supported division and its quantitative limit.}
Define, for $x\in[-L,L]$,
\[
 g_\lambda(x)=-\frac1\gamma
                  \int_{-L}^{x}\sin(\gamma(x-u))j_\lambda(u)du,
 \qquad g_\lambda(x)=0\quad(|x|>L).                \tag{SM.8}
\]
The conditions $\widehat j_\lambda(\pm\gamma)=0$
make both the value and first derivative at $L$
zero, as they already are at $-L$.
Thus this is a genuine compactly supported solution
of $(-\partial_x^2-\gamma^2)g_\lambda=j_\lambda$
without boundary delta terms.  It is even: reflecting
it gives another $L^2$ solution of the same equation,
and the difference is a linear combination of sine
and cosine, which cannot belong to $L^2(\mathbb R)$
unless it vanishes.
Fourier transformation gives the entire division
identity
\[
       \widehat g_\lambda(t)
           =\frac{\widehat j_\lambda(t)}{t^2-\gamma^2}.
                                                        \tag{SM.9}
\]
The double jets in (SM.5) make the quotient vanish
at both removed nodes.  This also gives a direct
Paley--Wiener verification: the singularities are
removable, the exponential type is still at most
$L$, and the quotient is square integrable on the
real line.  Formula (SM.8) supplies the support
assertion without relying solely on that theorem.
Since $j_\lambda\in H^s$ for $s<1/2$, weighted Plancherel gives
$g_\lambda\in H^{s+2}$ after the removable singularities
are filled in, so it belongs to the full
closed Weil form domain.

There is a nonzero limiting kernel
\[
 G(x)=-\frac1\gamma\int_{-\infty}^{x}
                        \sin(\gamma(x-u))K(u)du,
 \qquad \widehat G(t)=\frac{\Xi(t)}{t^2-\gamma^2}.
                                                        \tag{SM.10}
\]
Because $\widehat K(\pm\gamma)=0$, the same integral
can be written using the tail from $x$ to infinity
with the opposite sign.  AC-S04 therefore proves
rapid decay at both ends, indeed faster than every
exponential.  Its transform in (SM.10) is entire;
under the multiplicity hypothesis it still vanishes
at $\pm\gamma$.

For clarity, convergence of the compact inverse is
quantitative, despite the real Fourier denominator.
On $[-L,L]$ the sine integral in (SM.8) has $L^2$
operator norm at most $2L/\gamma$, by Young's
inequality.  Subtract (SM.10) and separate its
omitted left tail.  This gives
\[
 \|g_\lambda-G\|_2
 \leq C_\gamma L\|j_\lambda-K\|_{L^2(-L,L)}
   +C_\gamma\sqrt L\int_{-\infty}^{-L}|K(u)|du
   +\|1_{|x|>L}G\|_2\longrightarrow0.             \tag{SM.11}
\]
The first error is $O((\log\lambda)\lambda^{-1/2})$,
and the remaining tails decay faster than every
power of $\lambda$.
The limits $K,G$ are linearly independent: their
Fourier transforms differ by the nonconstant factor
$(t^2-\gamma^2)^{-1}$ on an open set where $\Xi\ne0$.
Consequently the Gram matrices of $j_\lambda,g_\lambda$
have a uniformly positive smallest eigenvalue for
large $\lambda$.

\proofstep{4. The two-dimensional variational estimate.}
By discreteness of the distinct zero ordinates,
\[
 d_\gamma=
   \inf_{\gamma'\ne\pm\gamma}
                            |(\gamma')^2-\gamma^2|>0.
\]
At the two excluded ordinates the transform of
$g_\lambda$ is exactly zero.
The full closed-domain sampling identity AC-S03
and (SM.9) therefore give, including every multiplicity,
\[
               q_L(g_\lambda)\leq d_\gamma^{-2}q_L(j_\lambda).
                                                        \tag{SM.12}
\]
Cauchy--Schwarz for the positive form bounds its
cross term on their span.  Together with the lower
Gram bound and (SM.6), this bounds the largest
Rayleigh quotient on that two-dimensional even
space by the first right side of (SM.1).
The minimum principle proves (SM.1), first in the
even sector and then in the unrestricted space.
The classical scalar leakage asymptotic supplies
its displayed exponential version.

If (SE.14) held and a multiple zero existed,
choose $\varepsilon<1+\alpha$ in (SM.1).
Its upper bound contradicts the lower bound by
a power of $\lambda$.  Hence (SE.14) would imply
that every zero is simple.  The argument asserts
only this conditional obstruction to that particular
sufficient estimate.
\end{proof}

\paragraph{What the two zero conditions mean.}
The Weil form samples values with multiplicity weights;
its null condition does not impose vanishing jets.
The full theta image has the stronger zero jets inherited
from its factor $\zeta$.  AC-S34 uses that distinction
exactly, by removing one factor at a hypothetical multiple
zero while retaining every sampled value zero.
It does not establish a counterexample for the actual
zeta function, logical independence from $A$, or a
necessary simplicity hypothesis for every possible
ground-eigenvector convergence argument.

\begin{proposition}[AC-S35: a ground line cannot encode direct-sum multiplicity]
\label{AC-S35}
Let $m\geq2$ be an integer and consider the entire
completed function $\Xi_m(t)=\Xi(t)^m$, corresponding
to the $m$-fold direct sum of the trivial Artin factor,
with the pole-removing completion taken componentwise.
Its Weil form and associated finite-window operator satisfy
\[
           q_L^{(m)}=m q_L,\qquad A_L^{(m)}=m A_L.
                                                        \tag{SM.13}
\]
Their lowest eigenspaces, including their multiplicities,
are identical for every $L$.  But the inverse Fourier
targets are $K$ and its $m$-fold convolution $K^{*m}$,
and these functions are not proportional.
Consequently a universal assertion of an eventually
simple lowest line whose scalar-normalized Fourier
transform recovers the entire completed function
cannot hold simultaneously for these two inputs.
This observation concerns such a universal assertion;
it does not refute the specific zeta conjecture in S25.
\end{proposition}
\begin{proof}
Every zero multiplicity is multiplied by $m$ under
the power operation.  The gamma and polar terms
of the componentwise completion are multiplied
by the same factor in the logarithmic explicit
formula.  Thus its Weil distribution and its
finite-window form are exactly $m$ times the
original ones; under $A$ this is also immediate
from the sampling formula AC-S03.
Multiplication of a closed lower semibounded form
by a positive scalar preserves its domain and
multiplies its associated self-adjoint operator
by that scalar, which proves (SM.13).

The function $K$ lies in $L^1\cap L^2$, by
AC-S04.  Repeated convolution and Young's
inequality therefore give $K^{*m}\in L^1\cap L^2$,
with Fourier transform $\Xi^m$.
If $K^{*m}=cK$, Fourier transformation would give
$\Xi(t)^m=c\Xi(t)$.
On an open set where $\Xi\ne0$, this would make
$\Xi^{m-1}$ constant, and hence $\Xi$ constant
by analyticity, a contradiction.

To include arbitrary scalar normalizations in the
asserted obstruction, suppose the common eventual
simple ground line had transforms $F_L$, and both
$c_LF_L\to\Xi$ and $d_LF_L\to\Xi^m$ locally uniformly.
Choose $t_0$ with $\Xi(t_0)\ne0$.
Evaluation at $t_0$ forces
$d_L/c_L\to\Xi(t_0)^{m-1}$.
The two limits would then be proportional everywhere,
contrary to the preceding argument.
Thus at least one of those simultaneous assertions
fails.  No choice is made here as to which one,
and no assertion of a multiple zero of $\zeta$
or independence from $A$ is used.
\end{proof}

\paragraph{Direct sums as a precise geometric dictionary.}
The categorical input in AC-S35 is elementary and
does not posit a new arithmetic surface.
For finite-dimensional Frobenius representations,
$\det(1-T(F\oplus G))=\det(1-TF)\det(1-TG)$,
so their local $L$-factors multiply under direct sum.
In the trivial Artin case, $m$ disjoint copies of
$\operatorname{Spec}\mathbb Q$ give the direct sum
of $m$ trivial representations, with finite Euler
product $\zeta(s)^m$.
The Weil trace expression is additive, whereas
the completed determinant is multiplicative.
A single ground line is unchanged by this direct-sum
scaling.  A realization intended to retain these
multiplicities must therefore keep additional
representation data, as the explicit Hilbert direct
sum in AC-S08 does.  This dictionary adds no
Grothendieck archive locator or unproved motivic
identification to the analytic construction.
